% !TEX program = xelatex
\documentclass[11pt,a4paper]{article}

\usepackage[UTF8]{ctex}
\usepackage[margin=2.2cm]{geometry}
\usepackage{amsmath, amssymb, amsthm}
\usepackage{listings}
\usepackage{xcolor}
\usepackage{tikz}
\usepackage{booktabs}
\usepackage{multirow}
\usepackage{hyperref}
\usetikzlibrary{positioning, arrows.meta, shapes.geometric, fit, calc, decorations.pathreplacing}

\hypersetup{
  colorlinks=true,
  linkcolor=blue!60!black,
  urlcolor=teal!70!black,
}

\lstdefinelanguage{Rust}{
  morekeywords={fn,let,mut,pub,struct,impl,use,mod,self,Self,Option,Some,None,return,if,else,for,while,loop,match,break,continue,as,where,trait,type,enum,move,ref,const,static,unsafe,extern,crate,super,dyn,async,await,box,in},
  sensitive=true,
  morecomment=[l]{//},
  morecomment=[s]{/*}{*/},
  morestring=[b]",
}

\lstset{
  language=Rust,
  basicstyle=\ttfamily\footnotesize,
  keywordstyle=\color{blue!60!black}\bfseries,
  commentstyle=\color{green!40!black}\itshape,
  stringstyle=\color{red!60!black},
  numbers=left,
  numberstyle=\tiny\color{gray},
  numbersep=8pt,
  frame=single,
  framerule=0.4pt,
  rulecolor=\color{gray!40},
  breaklines=true,
  breakatwhitespace=true,
  showstringspaces=false,
  tabsize=4,
  columns=fullflexible,
  keepspaces=true,
}

\title{\textbf{halfedge\,: 网格修复与洞填充模块设计文档} \\ \large \texttt{src/repair.rs}}
\author{halfedge 库}
\date{2026-07-05}

\begin{document}
\maketitle
\tableofcontents

\section{模块定位}

\texttt{repair.rs} 提供网格修复操作：洞填充、面删除、孤立元素清理、
非流形检测和一键修复 pipeline。适用于处理导入模型中的拓扑缺陷——
缺失面、退化面、孤立顶点等。

\subsection{API 总览}

\begin{center}
\renewcommand{\arraystretch}{1.18}
\begin{tabular}{@{}lp{3.8cm}l@{}}
  \toprule
  \textbf{函数} & \textbf{功能} & \textbf{返回值} \\
  \midrule
  \texttt{remove\_face} & 删除面，创建边界半边环 & \texttt{Result<(), TopErr>} \\
  \texttt{fill\_hole} & 填充单个洞 & \texttt{Result<Vec<FaceId>, TopErr>} \\
  \texttt{fill\_all\_holes} & 填充所有洞 & \texttt{Result<Vec<Vec<FaceId>>, TopErr>} \\
  \texttt{remove\_isolated\_vertices} & 删除孤立顶点 & \texttt{usize} \\
  \texttt{remove\_degenerate\_faces} & 删除退化面 & \texttt{usize} \\
  \texttt{detect\_nonmanifold\_edges} & 检测非流形边 & \texttt{Vec<HalfEdgeId>} \\
  \texttt{detect\_nonmanifold\_vertices} & 检测非流形顶点 & \texttt{Vec<VertexId>} \\
  \texttt{repair\_mesh} & 一键修复 pipeline & \texttt{Result<RepairStats, TopErr>} \\
  \bottomrule
\end{tabular}
\end{center}

\section{核心算法}

\subsection{remove\_face：删除面创建洞}

删除面 $f$ 及其关联的面侧半边，为每条暴露的边创建边界半边，
使洞边界形成完整的半边环供后续 \texttt{fill\_hole} 使用。

\subsubsection{流程图}

\begin{center}
\resizebox{\textwidth}{!}{%
\begin{tikzpicture}[
  node distance=0.45cm,
  proc/.style={draw, rounded corners=2pt, minimum width=11cm, minimum height=0.65cm, align=center, font=\small},
  dec/.style={diamond, draw, aspect=2.4, fill=orange!12, inner sep=1pt, font=\small, align=center, minimum width=3.5cm},
  op/.style={-Stealth, thick},
  startend/.style={draw, rounded corners=10pt, minimum width=2.6cm, minimum height=0.65cm, font=\small, fill=gray!12}
]
  \node[startend] (s) {开始：输入 \texttt{face}};
  \node[dec, below=of s] (d0) {面半边为空？};
  \node[proc, right=1.6cm of d0, fill=green!12] (early) {删除面，返回 \texttt{Ok(())}};
  \node[proc, below=of d0, fill=blue!8] (p1) {1. 收集面半边数据（克隆）};
  \node[proc, below=of p1, fill=blue!8] (p2) {2. 清空各 twin 的 twin 指向（$\text{twin.twin} \gets \text{None}$）};
  \node[proc, below=of p2, fill=blue!8] (p3) {3. 删除面和面侧半边};
  \node[dec, below=of p3] (d1) {twin 是边界半边？};
  \node[proc, right=1.6cm of d1, fill=red!12] (del) {删除该边界半边（边两侧均无面）};
  \node[proc, below=of d1, fill=yellow!15] (p4) {4a. 创建新边界半边 $b$（$\text{vertex} = \text{dst}$）};
  \node[proc, below=of p4, fill=yellow!15] (p5) {4b. 配对 $b \leftrightarrow \text{twin}$（$\text{next/prev}$ 保持 \texttt{None}）};
  \node[proc, below=of p5, fill=blue!8] (p6) {5. 修复所有顶点入口};
  \node[startend, below=of p6, fill=green!12] (e) {返回 \texttt{Ok(())}};
  \draw[op] (s) -- (d0);
  \draw[op] (d0) -- node[above, font=\scriptsize] {是} (early);
  \draw[op] (d0) -- node[right, font=\scriptsize] {否} (p1);
  \draw[op] (p1) -- (p2);
  \draw[op] (p2) -- (p3);
  \draw[op] (p3) -- (d1);
  \draw[op] (d1) -- node[above, font=\scriptsize] {是（相邻面已删）} (del);
  \draw[op] (d1) -- node[right, font=\scriptsize] {否} (p4);
  \draw[op] (p4) -- (p5);
  \draw[op] (p5) -- (p6);
  \draw[op] (p6) -- (e);
  \draw[op] (del) |- (p6);
\end{tikzpicture}
}
\end{center}

\subsubsection{边界半边的方向与 next/prev}

新边界半边 $b_i$ 与被删除的面半边 $h_i$ 同向：
\[
  b_i.\text{vertex} = h_i.\text{vertex} = \text{dst}_i
\]
边界半边\textbf{不设置} \texttt{next/prev}（保持 \texttt{None}），因为：
\begin{itemize}
  \item \texttt{BoundaryLoop} 通过顶点环遍历（\texttt{next\_boundary\_halfedge}），
        不依赖边界半边的 \texttt{next/prev}；
  \item 当相邻面被连续删除时，不维护 \texttt{next/prev} 可避免复杂的链表拼接。
\end{itemize}

\subsubsection{相邻面已删除的处理}

当删除面 $f_2$ 的某条面半边的 twin 是边界半边 $b$（说明相邻面 $f_1$ 已被删除），
直接删除 $b$ 而非创建新边界半边——此时该边两侧均无面，不存在「洞边界」。

\subsection{fill\_hole：洞填充}

填充由 \texttt{boundary\_he} 指定的单个洞，返回新创建的面 ID 列表。

\subsubsection{核心设计：不提前删除边界半边}

\texttt{fill\_hole} \textbf{不提前删除}边界半边，而是让
\texttt{add\_triangle} 的 twin 配对逻辑自然找到并删除它们。
原因：\texttt{add\_triangle} 的 twin 配对搜索条件为
$E.\text{vertex}=\text{src} \wedge E.\text{twin}.\text{vertex}=\text{dst}
\wedge E.\text{twin}.\text{face}=\text{None}$。
若提前删除边界半边，相邻面半边的 \texttt{twin} 变为 \texttt{None}，
配对失败。

\subsubsection{流程图}

\begin{center}
\resizebox{\textwidth}{!}{%
\begin{tikzpicture}[
  node distance=0.45cm,
  proc/.style={draw, rounded corners=2pt, minimum width=11cm, minimum height=0.65cm, align=center, font=\small},
  dec/.style={diamond, draw, aspect=2.4, fill=orange!12, inner sep=1pt, font=\small, align=center, minimum width=3.5cm},
  op/.style={-Stealth, thick},
  startend/.style={draw, rounded corners=10pt, minimum width=2.6cm, minimum height=0.65cm, font=\small, fill=gray!12}
]
  \node[startend] (s) {开始：输入 \texttt{boundary\_he}};
  \node[dec, below=of s] (d0) {\texttt{boundary\_he.face} $=$ \texttt{None}？};
  \node[proc, right=1.6cm of d0, fill=red!12] (err1) {返回 \texttt{Err(Inconsistent)}};
  \node[proc, below=of d0, fill=blue!8] (p1) {1. \texttt{BoundaryLoop} 收集边界环半边};
  \node[dec, below=of p1] (d1) {$n < 3$？};
  \node[proc, right=1.6cm of d1, fill=red!12] (err2) {返回 \texttt{Err(Inconsistent)}};
  \node[proc, below=of d1, fill=blue!8] (p2) {2. 提取顶点坐标};
  \node[proc, below=of p2, fill=yellow!15] (p3) {3. \texttt{ear\_clipping\_3d} 三角化};
  \node[proc, below=of p3, fill=yellow!15] (p4) {4. 清除边界半边 next/prev $\gets$ \texttt{None}};
  \node[proc, below=of p4, fill=blue!8] (p5) {5. 对每个三角形调用 \texttt{add\_triangle}};
  \node[proc, below=of p5, fill=green!12] (p6) {6. 修复顶点入口};
  \node[startend, below=of p6, fill=green!12] (e) {返回 \texttt{Ok(filled\_faces)}};
  \draw[op] (s) -- (d0);
  \draw[op] (d0) -- node[above, font=\scriptsize] {否} (err1);
  \draw[op] (d0) -- node[right, font=\scriptsize] {是} (p1);
  \draw[op] (p1) -- (d1);
  \draw[op] (d1) -- node[above, font=\scriptsize] {是} (err2);
  \draw[op] (d1) -- node[right, font=\scriptsize] {否} (p2);
  \draw[op] (p2) -- (p3);
  \draw[op] (p3) -- (p4);
  \draw[op] (p4) -- (p5);
  \draw[op] (p5) -- (p6);
  \draw[op] (p6) -- (e);
\end{tikzpicture}
}
\end{center}

\subsubsection{清除 next/prev 的必要性}

步骤 4 将边界半边的 \texttt{next/prev} 设为 \texttt{None}。原因：
当 \texttt{add\_triangle} 删除某个边界半边 $b_i$ 时，环中其他边界半边
$b_{i-1}$、$b_{i+1}$ 的 \texttt{next/prev} 可能仍指向 $b_i$，
造成悬空指针。清除后 \texttt{next = None}，\texttt{validate\_mesh}
跳过该项检查，不报错。

\subsubsection{twin 配对追踪（以三角形洞为例）}

设洞边界有 3 条半边 $b_0(A \to B)$、$b_1(B \to C)$、$b_2(C \to A)$，
其 twin 为面半边 $f_0$、$f_1$、$f_2$。
\texttt{BoundaryLoop} 收集 tip 序列 $[B, C, A]$，
\texttt{ear\_clipping\_3d} 返回 $[[0, 1, 2]]$，即调用
$\texttt{add\_triangle}(B, C, A)$：

\begin{center}
\renewcommand{\arraystretch}{1.18}
\begin{tabular}{@{}llll@{}}
  \toprule
  \textbf{新边} & \textbf{搜索条件} & \textbf{匹配} & \textbf{操作} \\
  \midrule
  $h_0\!: B \to C$ & $E.\text{vertex}=B$, $E.\text{twin}.\text{vertex}=C$ & $E = f_1$ & 删除 $b_1$，配对 $h_0 \leftrightarrow f_1$ \\
  $h_1\!: C \to A$ & $E.\text{vertex}=C$, $E.\text{twin}.\text{vertex}=A$ & $E = f_2$ & 删除 $b_2$，配对 $h_1 \leftrightarrow f_2$ \\
  $h_2\!: A \to B$ & $E.\text{vertex}=A$, $E.\text{twin}.\text{vertex}=B$ & $E = f_0$ & 删除 $b_0$，配对 $h_2 \leftrightarrow f_0$ \\
  \bottomrule
\end{tabular}
\end{center}

三条边界半边全部被消耗，洞完全填充。

\subsubsection{多三角形洞的内边处理}

对 $n > 3$ 的洞，\texttt{ear\_clipping\_3d} 产生多个三角形。
首三角形的非边界边（内边）由 \texttt{add\_triangle} 创建临时边界半边，
后续三角形的 \texttt{add\_triangle} 调用会找到该临时边界半边并配对，
最终所有边界半边被消耗。

\subsection{fill\_all\_holes}

枚举所有边界环，逐个调用 \texttt{fill\_hole}：

\begin{lstlisting}
let loops = boundary_loops(mesh);
for boundary in &loops {
    fill_hole(mesh, boundary[0])?;
}
\end{lstlisting}

\subsection{remove\_degenerate\_faces}

用 Shewchuk 鲁棒谓词 \texttt{is\_triangle\_degenerate\_3d} 检测
三顶点共线的退化面，调用 \texttt{remove\_face} 删除之。

\subsubsection{数据丢弃警告}

面删除失败时不再静默，会通过 \texttt{eprintln!} 输出警告：
\begin{center}
\texttt{[halfedge::remove\_degenerate\_faces] 警告：N 个退化面删除失败}
\end{center}

\subsection{非流形检测}

\subsubsection{detect\_nonmanifold\_edges}

统计每条无向边被几个面共享。流形网格中每条边至多 2 个面；
若 $> 2$ 则为非流形边。

\subsubsection{detect\_nonmanifold\_vertices}

对每个顶点检测 outgoing 环是否闭合。内部顶点应闭合，
边界顶点应为开链。若 CCW 和 CW 两个方向均不闭合且
边界半边数 $> 2$，则为非流形顶点。

\subsection{repair\_mesh：一键修复 pipeline}

\begin{center}
\resizebox{0.85\textwidth}{!}{%
\begin{tikzpicture}[
  node distance=0.6cm,
  proc/.style={draw, rounded corners=2pt, minimum width=8cm, minimum height=0.65cm, align=center, font=\small},
  op/.style={-Stealth, thick},
  startend/.style={draw, rounded corners=10pt, minimum width=2.6cm, minimum height=0.65cm, font=\small, fill=gray!12}
]
  \node[startend] (s) {输入：待修复网格};
  \node[proc, below=of s, fill=blue!8] (p1) {1. \texttt{fill\_all\_holes}：填充所有洞};
  \node[proc, below=of p1, fill=yellow!15] (p2) {2. \texttt{remove\_degenerate\_faces}：删除退化面};
  \node[proc, below=of p2, fill=yellow!15] (p3) {3. \texttt{remove\_isolated\_vertices}：删除孤立顶点};
  \node[proc, below=of p3, fill=green!12] (p4) {4. \texttt{fix\_orientations}：法线一致化};
  \node[startend, below=of p4, fill=green!12] (e) {返回 \texttt{RepairStats}};
  \draw[op] (s) -- (p1);
  \draw[op] (p1) -- (p2);
  \draw[op] (p2) -- (p3);
  \draw[op] (p3) -- (p4);
  \draw[op] (p4) -- (e);
\end{tikzpicture}
}
\end{center}

\texttt{RepairStats} 记录各项操作的统计信息：
\begin{center}
\renewcommand{\arraystretch}{1.18}
\begin{tabular}{@{}ll@{}}
  \toprule
  \textbf{字段} & \textbf{含义} \\
  \midrule
  \texttt{holes\_filled} & 填充的洞数量 \\
  \texttt{isolated\_vertices\_removed} & 删除的孤立顶点数 \\
  \texttt{degenerate\_faces\_removed} & 删除的退化面数 \\
  \texttt{faces\_flipped} & 法线一致化翻转的面数 \\
  \bottomrule
\end{tabular}
\end{center}

\section{复杂度}

\begin{center}
\renewcommand{\arraystretch}{1.18}
\begin{tabular}{@{}lll@{}}
  \toprule
  \textbf{操作} & \textbf{时间复杂度} & \textbf{备注} \\
  \midrule
  \texttt{remove\_face}              & $O(k + V)$   & $k$ = 面边数，$V$ = 顶点修复 \\
  \texttt{fill\_hole}                & $O(n^2 \cdot H)$ & $n$ = 洞顶点数，$H$ = 半边总数（twin 搜索） \\
  \texttt{fill\_all\_holes}          & $O(\sum n_i^2 \cdot H)$ & 各洞独立处理 \\
  \texttt{remove\_isolated\_vertices}& $O(V)$       & 线性扫描 \\
  \texttt{remove\_degenerate\_faces} & $O(F \cdot H)$ & $F$ = 面数，$H$ = 半边搜索 \\
  \texttt{detect\_nonmanifold\_edges}& $O(H)$       & 单遍扫描 \\
  \texttt{detect\_nonmanifold\_vertices} & $O(V \cdot \bar{d})$ & $\bar{d}$ = 平均度数 \\
  \texttt{repair\_mesh}              & 上述之和       & pipeline 顺序执行 \\
  \bottomrule
\end{tabular}
\end{center}

\section{退化守卫}

\begin{itemize}
  \item \texttt{fill\_hole}：校验 \texttt{boundary\_he.face = None}，
        否则返回 \texttt{Inconsistent}；
  \item \texttt{fill\_hole}：边界环顶点数 $< 3$ 时返回 \texttt{Inconsistent}；
  \item \texttt{remove\_face}：面半边为空时直接删除面并返回；
  \item \texttt{remove\_face}：相邻面已删除时（twin 是边界半边），直接删除
        该边界半边而非创建新边界；
  \item \texttt{remove\_degenerate\_faces}：使用 Shewchuk 鲁棒谓词，
        避免浮点阈值误判；
  \item \texttt{detect\_nonmanifold\_vertices}：使用 \texttt{halfedge\_count + 1}
        作为最大迭代次数，防止死循环；
  \item 所有顶点入口修复使用 \texttt{find\_outgoing\_halfedge} 确保引用有效。
\end{itemize}

\section{测试覆盖}

\begin{center}
\renewcommand{\arraystretch}{1.18}
\begin{tabular}{@{}ll@{}}
  \toprule
  \textbf{测试名} & \textbf{验证点} \\
  \midrule
  \texttt{fill\_hole\_basic} & 单面洞填充后无残留边界 \\
  \texttt{fill\_all\_holes\_icosphere} & icosphere 单面洞一键填充 \\
  \texttt{fill\_hole\_on\_closed\_mesh\_returns\_empty} & 闭合网格无洞 \\
  \texttt{fill\_hole\_non\_boundary\_is\_error} & 非边界半边报错 \\
  \texttt{remove\_isolated\_vertices\_basic} & 删除无关联顶点 \\
  \texttt{remove\_face\_creates\_boundary} & 删除面后产生 3 条边界半边 \\
  \texttt{remove\_face\_and\_fill\_restores} & 删除+填充恢复面数 \\
  \texttt{detect\_nonmanifold\_edges\_closed\_mesh} & 闭合无非流形边 \\
  \texttt{repair\_mesh\_basic} & 一键修复统计正确 \\
  \texttt{fill\_hole\_quad\_hole} & grid 边界洞填充 \\
  \texttt{remove\_and\_fill\_multiple\_faces} & 删除 2 面+填充 \\
  \texttt{remove\_degenerate\_faces\_on\_good\_mesh} & 好网格无退化面 \\
  \bottomrule
\end{tabular}
\end{center}

全模块共 12 个单元测试；项目整体 \texttt{cargo test} 共 450 个用例。

\section{设计权衡}

\begin{itemize}
  \item \textbf{不删除边界半边 vs 提前删除}：
        \texttt{fill\_hole} 保留边界半边让 \texttt{add\_triangle}
        自然配对，而非提前删除后手动重建 twin。后者需
        手动建立 twin 关系，逻辑更复杂且易遗漏。

  \item \textbf{边界半边 next/prev = None}：
        \texttt{BoundaryLoop} 用顶点环遍历，
        不依赖边界半边的 \texttt{next/prev}。设为
        \texttt{None} 避免连续删除相邻面时的链表拼接问题，
        \texttt{validate\_mesh} 对 \texttt{next = None} 跳过检查。

  \item \texttt{ear\_clipping\_3d} \textbf{而非扇形三角化}：
        对凹多边形洞，扇形三角化可能产生翻转面；
        \texttt{ear\_clipping\_3d} 投影到主平面后使用耳切法，
        能正确处理凹洞。

  \item \textbf{非流形检测 vs 修复}：
        当前仅提供检测，不自动修复非流形拓扑。
        修复策略取决于应用场景，过早修复可能破坏用户意图。
\end{itemize}

\end{document}
